\documentclass[letterpaper, 12pt]{article}
\usepackage[utf8]{inputenc}
\usepackage[T1]{fontenc}
\usepackage[spanish]{babel}
\usepackage{amsmath, amssymb, amsfonts}
\usepackage{geometry}
\usepackage{fancyhdr}
\usepackage{multicol}
\usepackage{xcolor}
\usepackage{tcolorbox}
\tcbuselibrary{skins}
\usepackage{enumitem}
\usepackage{graphicx}
\usepackage{hyperref}

% --- Configuración de Márgenes ---
\geometry{letterpaper, margin=1.5cm, top=3cm, bottom=2.5cm, headheight=2cm}

% --- Colores Usach Premium ---
\definecolor{celesteSuave}{HTML}{F0F7FF}
\definecolor{azulFuerte}{HTML}{1A5276}
\definecolor{azulClaro}{RGB}{120, 160, 210}

\pagestyle{fancy}
\fancyhf{} 

% --- ENCABEZADO ---
\newcommand{\logoup}{./images/logo_up.png}
\newcommand{\logousach}{./images/logo_usach.png}
\fancyhead[L]{\includegraphics[height=1.2cm]{\logoup}}
\fancyhead[R]{%
\begin{minipage}[b]{0.4\textwidth}
\raggedleft
\fontsize{9pt}{11pt}\selectfont
\textbf{\color{azulFuerte}Usach Premium} \\
Álgebra 2026 \\
\textit{\color{gray}@usach.premium}
\end{minipage} \hspace{0.1cm}
\includegraphics[height=1.2cm]{\logousach}
}

% --- PIE DE PÁGINA ---
\fancyfoot[L]{\fontsize{9pt}{11pt}\selectfont \textit{``Por y para estudiantes''}}
\fancyfoot[R]{\fontsize{9pt}{11pt}\selectfont Página \thepage}

\renewcommand{\headrulewidth}{0.4pt}
\renewcommand{\footrulewidth}{0.4pt}

% ============================================================
% CAJAS Y ENTORNOS REUTILIZABLES
% ============================================================
\newenvironment{solbox}[1]{%
    \begin{tcolorbox}[colback=celesteSuave, colframe=azulFuerte,
        title=\textbf{#1}, fonttitle=\bfseries]
}{%
    \end{tcolorbox}
}

\newenvironment{ejercicios}{}{}
\newenvironment{soluciones}{%
    \newpage
    \begin{center}
        \begin{tcolorbox}[colback=azulFuerte, colframe=azulFuerte,
            colupper=white, width=0.8\textwidth, center, arc=10pt]
            \centering\Large\textbf{SOLUCIONARIO}
        \end{tcolorbox}
    \end{center}
    \footnotesize
}{}

\begin{document}

% ==========================================
% PORTADA
% ==========================================
\vspace*{0.5cm} 

\begin{center}
\color{azulFuerte}
{\huge \textbf{GUÍA DE EJERCICIOS PEP 2}} \\
\vspace{0.5cm}
{\Large \textbf{Guía 7 - Matrices, Sumatorias y Complejos}} \\
\vspace{0.5cm}
{\Huge \textbf{Álgebra 1}} \\
\vspace{1.5cm}
\begin{tcolorbox}[colback=celesteSuave, colframe=azulFuerte, arc=10pt, width=0.85\textwidth, center, boxrule=1.5pt]
\centering
\vspace{0.2cm}
{\Large \textbf{Contenidos de la Guía}}
\vspace{0.3cm}
\color{black}
\begin{itemize}[leftmargin=3cm]
\item[\color{azulFuerte} $\Rightarrow$] Matrices, Determinantes y Sistemas de Ecuaciones.
\item[\color{azulFuerte} $\Rightarrow$] Sumatorias.
\item[\color{azulFuerte} $\Rightarrow$] Números Complejos.
\end{itemize}
\vspace{0.2cm}
\end{tcolorbox}
\end{center}

\vspace{0.5cm}
\begin{center}
{\Large \textbf{1er. Semestre de 2026}} \\
\vspace{1.3cm}
{\large \textbf{Usach Premium}} \\
\vspace{1.5cm}
\textit{``El fracaso no consiste en fallar, sino en no intentarlo.''}
\end{center}

\newpage
\pagenumbering{arabic}
\thispagestyle{fancy}

% ============================================================
% EJERCICIOS
% ============================================================
\begin{ejercicios}

\section*{I. Matrices y Sistemas de Ecuaciones}

1) Resuelva los siguientes sistemas de ecuaciones lineales:
$$ \text{i) } \begin{cases} x + 3y - z = -3 \\ 3x - y + 2z = 1 \\ 2x - y + z = -1 \end{cases} \qquad \text{ii) } \begin{cases} x + 3y - 3z = -5 \\ 2x - y + z = -3 \\ -6x + 3y - 3z = 4 \end{cases} \qquad \text{iii) } \begin{cases} 2x - y + z = 0 \\ x - y - 2z = 0 \\ 2x - 3y - z = 0 \end{cases} $$

\vspace{0.5cm}

2) Sean las matrices:
$$ A = \begin{pmatrix} 1 & 2 \\ 0 & -3 \end{pmatrix}, \quad B = \begin{pmatrix} 2 & -1 \\ 3 & 1 \end{pmatrix}, \quad C = \begin{pmatrix} 3 & 1 \\ -2 & 0 \end{pmatrix} $$
Compruebe el enunciado:
$$ \text{a) } (A+B)(A-B) \neq A^2 - B^2, \quad \text{donde } A^2 = AA \text{ y } B^2 = BB. $$

\vspace{0.5cm}

3) Dadas las siguientes matrices:
Calcule $A+B$ cuando la operación esté definida:
\begin{enumerate}[label=\alph*), leftmargin=1.5cm, itemsep=0.8em]
    \item[\textbf{b)}] $A = \begin{pmatrix} 0 & -2 & 7 \\ 5 & 4 & -3 \end{pmatrix}, \qquad B = \begin{pmatrix} 8 & 4 & 0 \\ 0 & 1 & 4 \end{pmatrix}$
    \item[\textbf{c)}] $A = \begin{pmatrix} 7 \\ -16 \end{pmatrix}, \qquad B = \begin{pmatrix} -11 \\ 9 \end{pmatrix}$
    \item[\textbf{d)}] $A = \begin{pmatrix} 2 & 1 \end{pmatrix}, \qquad B = \begin{pmatrix} 3 & -1 & 5 \end{pmatrix}$
\end{enumerate}

\vspace{0.5cm}

4) Sean las matrices $A, B$ definidas por:
$$ A = \begin{pmatrix} 1 & 1 & 0 \\ 0 & 1 & 1 \\ 1 & 0 & -1 \end{pmatrix}, \qquad B = \begin{pmatrix} x & 1 & -1 \\ 1 & x & 0 \\ 2 & -2 & 1 \end{pmatrix} $$
Resuelva la siguiente ecuación:
$$ |A^2| = |2B| $$

\vspace{0.5cm}

\newpage
\thispagestyle{fancy}
5) Sea la matriz $A$ definida por:
$$ A = \begin{pmatrix} 0 & 1 & -1 \\ 1 & 1 & 0 \\ 0 & 1 & 0 \end{pmatrix} $$
Resuelva la siguiente ecuación:
$$ |2A - xI_3| = 2x^2 - x^3 $$

\vspace{0.5cm}

6) Usando la regla de Cramer, encuentre el valor de $x \in \mathbb{R}$ que sea solución del sistema:
$$ \begin{cases} 2x - 3y + z = -2 \\ 3x - 2z = 4 \\ 2y - 4z = -3 \end{cases} $$

\vspace{0.5cm}

7) Dado el sistema de ecuaciones lineales de orden $3 \times 3$:
$$ \begin{cases} x - y = 1 \\ -x + (a + 4)y + 2z = a + 2 \quad (*) \\ 2x - 2y + (a - 2)z = 4 - a \end{cases} $$
\begin{enumerate}[label=(\alph*), leftmargin=1.2cm, itemsep=0.5em]
    \item Determine el conjunto:
    $$ S = \{a \in \mathbb{R} : (*) \text{ tenga solución única}\} $$
    \item Pruebe, usando la regla de Cramer, que $z = -1$ para todo $a \in S$.
\end{enumerate}

\vspace{0.5cm}

\newpage
\thispagestyle{fancy}
8) Sea la matriz $A \in M_3(\mathbb{R})$ definida por:
$$ A = \begin{pmatrix} -7 & 4 & 4 \\ 4 & -1 & 8 \\ 4 & 8 & -1 \end{pmatrix} $$
Determine el conjunto:
$$ S = \left\{ x \in \mathbb{R} : \left(\frac{x^2 - 5}{9}\right) A^2 - 3(2x^2 + 1)I_3 = (0)_3 \right\} $$

\vspace{0.5cm}

9) Sean las matrices $A, B \in M_3(\mathbb{R})$ y $\lambda \in \mathbb{R}$ definidas por:
$$ A = \begin{pmatrix} \lambda & 1 & 2 \\ 0 & 2 & -1 \\ 3 & 0 & \lambda \end{pmatrix}, \qquad B = \begin{pmatrix} 1 & 1 & 0 \\ 1 & 0 & -2 \\ 0 & 0 & 1 \end{pmatrix} $$
Determine el conjunto:
$$ S = \{ \lambda \in \mathbb{R} : |A| = |A - \lambda B^2| \} $$

\vspace{1cm}

\newpage
\thispagestyle{fancy}
\section*{II. Sumatorias}

1) Calcular la siguiente expresión:
$$ S = \sum_{\ell=2}^{n+2} (\ell - 1)^4 - \sum_{k=1}^n k^4 $$

\vspace{0.5cm}

2) Calcule el valor de la sumatoria:
$$ \sum_{i=1}^{17} \left[ \frac{i}{i^2 + 2i + 1} - \frac{i - 1}{i^2} \right] $$

\vspace{0.5cm}

3) Calcule la siguiente sumatoria:
$$ \sum_{i=2}^{15} \frac{3}{(2i - 1)(2i + 3)} $$

\vspace{0.5cm}

4) Sabiendo que $\displaystyle \sum_{k=1}^n a_k = \frac{n^2 + 5n}{2}$, calcule el valor de $a_5 + a_6$.

\vspace{0.5cm}

5) Encuentre el valor de $c \in \mathbb{N}_0$, si es que existe, que cumpla la siguiente igualdad:
$$ \sum_{k=2}^{c+1} \left( \frac{3}{4k^2 - 1} \right) = \frac{c}{17} $$

\vspace{0.5cm}

6) Demuestre que para todo $n \in \mathbb{N}$:
$$ \sum_{k=1}^n (2k - 1)^2 = \frac{n(4n^2 - 1)}{3} $$

\vspace{0.5cm}

7) Sabiendo que
$\displaystyle \sum_{k=2}^{12}a_k^2=432$ y
$\displaystyle \sum_{k=2}^{12}(a_k-3)^2=396$, determine el valor de:
$$ \sum_{k=2}^{12} (a_k^2 - 6a_k) $$

\vspace{1cm}

\section*{III. Números Complejos}

1) Calcule y exprese en la forma $a+bi$ el siguiente número complejo:
$$ z = \frac{(1+i)^6}{(1-i)^4} $$
 
\vspace{0.5cm}
 
2) Determine el lugar geométrico de los números complejos $z \in \mathbb{C}$ que satisfacen la siguiente condición, expresando el resultado como una ecuación en $x$ e $y$ (con $z = x+iy$):
$$ |z - 2i| = |z + 4| $$
 
\vspace{0.5cm}
 
3) Resuelva en $\mathbb{C}$ la siguiente ecuación, expresando cada solución en la forma $a+bi$:
$$ z^2 = 3 - 4i $$
 
\vspace{0.5cm}
 
4) Sean los números complejos $z_1 = 1 + i\sqrt{3}$ y $z_2 = \sqrt{3} - i$. Usando forma polar y el Teorema de De Moivre, calcule el siguiente producto y exprese el resultado en la forma $a+bi$:
$$ z_1^5 \cdot \overline{z_2}^{\,3} $$
 
\vspace{0.5cm}
 
5) Sea el polinomio $P(z) = z^3 + az^2 + bz + c$, con $a, b, c \in \mathbb{R}$. Determine los valores de $a$, $b$ y $c$ para que $P(z)$ sea divisible por el polinomio:
$$ Q(z) = z^2 - iz - (1+i) $$

\vspace{0.5cm}

6) Determine el valor de $k \in \mathbb{R}$ para que el número complejo
$$ z = \frac{2 + ki}{3 - i} $$
tenga módulo igual a $\sqrt{2}$.

\vspace{0.5cm}

7) Sea $z \in \mathbb{C}$ tal que $|z| = 1$ y $z \neq -1$. Demuestre que el número complejo
$$ w = \frac{z - 1}{z + 1} $$
es imaginario puro. $\star$


\newpage

% ============================================================
% SOLUCIONARIO
% ============================================================
\begin{soluciones}

\begin{solbox}{I. Respuestas - Matrices y Sistemas de Ecuaciones}
\begin{enumerate}[label=\color{azulFuerte}\bfseries\arabic*), leftmargin=*, itemsep=0.6em]
    \item \textbf{i)} $(x, y, z) = (-2, 1, 4)$ \quad \textbf{ii)} Sistema Incompatible (sin solución) \quad \textbf{iii)} $(x, y, z) = (0, 0, 0)$

    \item $(A+B)(A-B) = \begin{pmatrix} -6 & 5 \\ 3 & 17 \end{pmatrix} \quad \neq \quad A^2 - B^2 = \begin{pmatrix} 0 & -1 \\ -9 & 11 \end{pmatrix}$

    \item \textbf{b)} $A + B = \begin{pmatrix} 8 & 2 & 7 \\ 5 & 5 & 1 \end{pmatrix} \qquad$ \textbf{c)} $A + B = \begin{pmatrix} -4 \\ -7 \end{pmatrix} \qquad$ \textbf{d)} No definida.

    \item $\mathbf{x = -1}$

    \item $\mathbf{x = 2}$

    \item $\mathbf{x = 3}$

    \item \textbf{(a)} $\mathbf{S = \mathbb{R} \setminus \{-3, 2\}} \qquad$ \textbf{(b)} Demostrado ($z = -1$ para todo $a \in S$).

    \item $\mathbf{S = \{-4, 4\}}$

    \item $\mathbf{S = \left\{0, \frac{3}{2}\right\}}$
\end{enumerate}
\end{solbox}

\vspace{0.5cm}

\begin{solbox}{II. Respuestas - Sumatorias}
\begin{enumerate}[label=\color{azulFuerte}\bfseries\arabic*), leftmargin=*, itemsep=0.6em]
    \item $S = (n+1)^4 = \mathbf{n^4 + 4n^3 + 6n^2 + 4n + 1}$

    \item $\mathbf{\frac{17}{324}}$

    \item $\mathbf{\frac{602}{1705}}$

    \item $a_5 + a_6 = \mathbf{15}$

    \item $\mathbf{c = 0 \quad \text{o} \quad c = 7}$

    \item Demostrado (se verifica por inducción matemática para todo $n \in \mathbb{N}$).

    \item $\mathbf{297}$
\end{enumerate}
\end{solbox}

\vspace{0.5cm}

\begin{solbox}{III. Respuestas - Números Complejos}
\begin{enumerate}[label=\color{azulFuerte}\bfseries\arabic*), leftmargin=*, itemsep=0.6em]
    \item $\mathbf{z = 2i}$
 
    \item $\mathbf{y = -2x - 3}$ (recta en el plano complejo)
 
    \item $\mathbf{z = 2 - i} \quad \text{o} \quad \mathbf{z = -2 + i}$
 
    \item $\mathbf{128\sqrt{3} + 128i}$
 
    \item $\mathbf{a=-1,\ b=0,\ c=2}$. En efecto,
    $P(z)=Q(z)(z-1+i)$.
    
    \item $\mathbf{k = 4} \quad \text{o} \quad \mathbf{k = -4}$

    \item Demostrado: como $|z|=1 \Rightarrow \bar z = \dfrac{1}{z}$, se obtiene $\bar w = -w$, por lo tanto $\text{Re}(w) = 0$ y $w$ es imaginario puro.

\end{enumerate}
\end{solbox}
 
\vspace{0.8cm}
\begin{center}
    \small\textbf{Usach Premium} — ¡Nos vemos en la ayudantía! \\
    \textit{``Por y para estudiantes.''}
\end{center}
 
\end{soluciones}
 
\end{ejercicios}
\end{document}
